\section{Supplementary proofs and definitions}
\label{sec:suppl-definitions}

\subsection{Proofs}
\label{ssec:suppl-proofs}

Section~\ref{ssec:topogen-limitations} states that for any closed orientable 2-manifold, the first and second Betti numbers satisfy $\beta_1 = 2g$ and $\beta_2 = \beta_0$. We provide a proof of this below.

\noindent
A closed orientable surface of genus $g$ can be seen as a sphere with $g$ handles attached. 
For such a surface, the Euler characteristic is known to be
\begin{equation}
\chi = 2 - 2g.
\end{equation}

The Euler characteristic also satisfies
\begin{equation}
\chi = \beta_0 - \beta_1 + \beta_2,
\end{equation}
where $\beta_i$ are the Betti numbers.

Since the surface is closed and connected, we have $\beta_0 = 1$. 
For orientable surfaces, there is one independent 2-dimensional hole, so $\beta_2 = 1$.

Substituting these values gives
\begin{equation}
2 - 2g = 1 - \beta_1 + 1 \quad \Rightarrow \quad \beta_1 = 2g.
\end{equation}

Thus,
\begin{equation}
\beta_1 = 2g, \qquad \beta_2 = \beta_0.
\end{equation}


% \subsection{Metrics}

% \paragraph{Balanced accuracy} For a classification task with $K$ classes, the balanced accuracy is defined as:
% \begin{equation}\label{eq:balanced-accuracy}
% \text{Balanced Accuracy} = \frac{1}{K} \sum_{i=1}^{K} \frac{TP_i}{TP_i + FN_i}
% \end{equation}

% where $TP_i$ and $FN_i$ are the number of true positives and false negatives for class $i$, respectively. This metric is particularly useful for imbalanced datasets, as each class is reweighted according to its importance.